\documentclass[11pt,a4paper]{article}

% --- Packages ---
\usepackage[utf8]{inputenc}
\usepackage[T1]{fontenc}
\usepackage{amsmath,amssymb,amsthm,mathtools}
\usepackage{booktabs}
\usepackage{microtype}
\usepackage{hyperref}
\usepackage{cleveref}
\usepackage{enumitem}
\usepackage{geometry}
\usepackage{xcolor}

\geometry{
  a4paper,
  top=2.5cm,
  bottom=2.5cm,
  left=2.5cm,
  right=2.5cm
}

\hypersetup{
  colorlinks=true,
  linkcolor=blue!70!black,
  citecolor=green!50!black,
  urlcolor=blue!80!black
}

% --- Theorem Environments ---
\theoremstyle{definition}
\newtheorem{checklist}{Checklist}

\theoremstyle{remark}
\newtheorem{remark}{Remark}

% --- Math Shortcuts ---
\DeclareMathOperator{\RePart}{Re}
\DeclareMathOperator{\ImPart}{Im}
\renewcommand{\Re}{\RePart}
\renewcommand{\Im}{\ImPart}

% --- Metadata ---
\title{\textbf{Kernel Verification is Necessary But Insufficient}\\[0.5em]
\large Vacuity Detection as the Dominant Failure Mode in LLM-Assisted Formalization}

\author{\textbf{Xavier Fresquet}\thanks{Sorbonne Center for Artificial Intelligence (SCAI), Sorbonne Universit\'e, Paris, France. Correspondence: \texttt{scai@sorbonne-universite.fr}}}

\date{August 2026}

\begin{document}

\maketitle

\begin{abstract}
We document a systematic study of the failure modes encountered in a large-scale LLM-assisted formalization of a frontier open problem (the Riemann Hypothesis). Machine-certified Lean~4 proofs with zero custom axioms do not guarantee mathematical soundness; instead, we identify three dominant failure modes that kernel verification cannot detect: (1) legacy false axioms surviving apparent proof success (e.g., $|M(x)| \le 2$, refuted by $M(33)=3$), (2) subtle machinery conflations (regulator confusion across Stages 4--5), and (3) vacuous impossibility theorems (formally correct but quantifying over zero-element sets).

We present a methodology for vacuity auditing applicable to any LLM-formalized mathematics and demonstrate how close reading, multi-model auditing, and explicit scope specification prevent these failures. One real architectural diagnostic emerges: the $q^2$ majorant loss in functional-equation machinery, which forces a switch from classical absolute bounds to signed cancellation strategies. This is a substantive frontier discovery, not an impossibility claim, and illustrates what sound negative results look like.

We conclude with recommendations for transparent LLM-assisted mathematical research: kernel verification is syntactic correctness, not semantic validity. Vacuity detection requires human interpretation of quantifier scope, cannot be automated, and must be built into formalization workflows as a mandatory audit gate.

\textbf{Keywords:} Formal Verification; Lean 4; LLM-Assisted Mathematics; Vacuity Testing; Computational Epistemology; Failure Analysis.
\end{abstract}

% =============================================================================
\section{Introduction}
\label{sec:intro}
% =============================================================================

Large Language Models have become standard tools in mathematical formalization, translating paper proofs into kernel-checked formal systems (Lean, Isabelle, Coq). Kernel verification ensures syntactic correctness: the type-checker confirms that all terms are well-formed and all proofs are logically valid.

But syntactic correctness is not semantic validity. A theorem that says ``for all $x$ in the empty set, $P(x)$'' is syntactically correct (vacuously true) and kernel-verified, yet tells us nothing.

This paper documents a case study where this gap became critical: a multi-stage LLM-assisted formalization of the Riemann Hypothesis that produced 8,649 kernel-verified jobs with zero custom axioms, yet contained three distinct classes of errors that kernel verification could not detect.

\subsection{The Central Problem}

When a machine-certified proof passes the Lean kernel, we know:
\begin{itemize}[leftmargin=2em]
  \item The syntax is correct (no type errors, all tactics apply)
  \item The logical structure is sound (no inference rules violated)
  \item The axiom inventory is transparent (we can list every assumption)
\end{itemize}

We do \emph{not} know:
\begin{itemize}[leftmargin=2em]
  \item Whether the hypothesis has non-empty models (vacuity)
  \item Whether the conclusion captures the intended problem (scope mismatch)
  \item Whether a supporting lemma was actually used, or if it's a dead weight (utility audit)
\end{itemize}

All three gaps appeared in our formalization. This paper is a post-mortem.

\subsection{Contributions}

\begin{enumerate}[label=\textbf{\arabic*.}]
  \item A taxonomy of three failure modes in LLM-assisted formalization with concrete examples
  \item A methodology for vacuity auditing, including a checklist applicable to any kernel-verified impossibility claim
  \item One real (non-vacuous) architectural diagnostic: the $q^2$ majorant loss that forces signature changes in proof strategy
  \item Recommendations for transparent LLM-assisted research workflows
\end{enumerate}

% =============================================================================
\section{The RH Formalization Experiment}
\label{sec:experiment}
% =============================================================================

\subsection{Setup}

Between August 2026 and the present, we conducted a large-scale formalization of the Riemann Hypothesis via:

\begin{itemize}[leftmargin=2em]
  \item 498 peer-reviewed papers, indexed in a knowledge graph (21,942 nodes, 12,859 relations)
  \item Multi-LLM parliament: Claude Sonnet 5, GPT-5.6 Sol, Gemini 3.1, Mistral Medium 3.5, DeepSeek-v4, Kimi K3
  \item Nine formal reduction stages from Nyman--Beurling criterion to the RH frontier
  \item Active formalization umbrella: \texttt{NBMellinTools} (47 modules, 290+ lemmas)
  \item Exploratory frontier family: NB15 modules (17 files, 4,645 lines)
\end{itemize}

\subsection{Apparent Success (Stage 9, Early August)}

By August 3, 2026, the formalization appeared complete:
\begin{itemize}[leftmargin=2em]
  \item 8,649 Lean jobs verified
  \item 0 custom axioms (only Lean/Mathlib foundations)
  \item 0 \texttt{sorry} statements
  \item Frontier condition (H15) precisely characterized
  \item Four ``formally proved barriers'' (impossibility results) documented
  \item Full audit trail and knowledge graph published
\end{itemize}

From a syntactic perspective, this was a success: the kernel certified every statement.

\subsection{The Hidden Reality (August 3--4)}

Upon adversarial auditing and close reading:

\begin{enumerate}[label=\textbf{\arabic*.}]
  \item Three legacy false axioms were discovered (should have been purged earlier)
  \item One critical machinery conflation in Stages 4--5 (regulator vs.~damping)
  \item Four claimed ``barriers'' were vacuous (quantified over zero-element sets or single hand-crafted examples)
  \item One real diagnostic emerged: the $q^2$ loss in functional-equation machinery
\end{enumerate}

The formalization was not wrong; it was incomplete in its self-assessment. The kernel checked syntax, but humans had to audit semantics.

% =============================================================================
\section{Failure Mode 1: Legacy False Axioms}
\label{sec:failure1}
% =============================================================================

\subsection{The Three Cases}

\begin{description}
  \item[\texttt{mobius\_partial\_sum\_bounded}:] Asserted $\forall x \in \mathbb{R}^+, |M(x)| \le 2$ where $M(x) = \sum_{n \le x} \mu(n)$ is the Möbius summatory function.
    \textbf{False:} $M(33) = 3$. (This is a well-known violation.)

  \item[\texttt{weil\_bound\_mobius}:] Attempted to apply Weil bounds from finite-field character sums ($\mathbb{F}_q[t]$) directly to Möbius oscillation on $\mathbb{Z}$.
    \textbf{Error:} Character bounds on finite fields do not transfer to integer oscillation without substantial reformulation.

  \item[\texttt{ZetaZero}:] Defined a vacuous predicate to represent ``all RH-level consequences'' without specifying what it contained.
    \textbf{Problem:} The predicate had zero inhabitants; any universal statement about it is vacuously true.
\end{description}

\subsection{Why They Survived}

Each axiom was stated as a hypothesis in an implication, never directly used to derive \texttt{False}. The structure was:

\begin{verbatim}
lemma some_theorem (h : mobius_partial_sum_bounded) : ... :=
  ...
\end{verbatim}

The kernel verified this: IF \texttt{mobius\_partial\_sum\_bounded} were true, THEN the conclusion follows. The conditional is syntactically valid regardless of whether the hypothesis is true.

\subsection{How They Were Caught}

Gemini 3.1 Pro, running with 1M-token context window on the full source tree, performed a forensic audit:
\begin{enumerate}[label=(\roman*)]
  \item Identified every \texttt{axiom} and \texttt{sorry} declaration
  \item Cross-checked each against published literature
  \item Flagged mismatches between formalized statements and cited sources
\end{enumerate}

Multi-model consensus (Claude, GPT, Gemini, Mistral) confirmed all three were erroneous.

\subsection{Lesson}

Kernel verification does not substitute for literature review. A theorem can be conditionally true ($A \to B$) while its hypothesis ($A$) is false in the real world.

For LLM-assisted formalization:
\begin{itemize}[leftmargin=2em]
  \item \textbf{Mandatory:} Cross-check every axiom against peer-reviewed sources before final publication
  \item \textbf{Mandatory:} When importing classical results, cite the original, not the translation
  \item \textbf{Useful:} Build axiom-checking into CI/CD (automated but not sufficient)
\end{itemize}

% =============================================================================
\section{Failure Mode 2: Machinery Conflation}
\label{sec:failure2}
% =============================================================================

\subsection{The Regulator Confusion (Stages 4--5)}

In Stages 4--5, the formalization attempted to represent the ``regulator'' (a key scaling factor in Dirichlet series machinery) as a bidisc damping term:
\[
\text{Attempted}: \quad \rho^m \sigma^\ell \quad (\text{bidisc damping})
\]

However, the correct Estermann regulator is:
\[
\text{Classical}: \quad \tau^{m\ell} = (d(n))^{m\ell} \quad (\text{divisor function})
\]

These are distinct structures. The bidisc damping is a coordinate scaling; the Estermann regulator is an arithmetic function.

\subsection{Why It Survived the Kernel}

The formalization with the wrong regulator was still \emph{algebraically correct}. The functional-equation manipulations remained valid because the equations only depend on the regulator's scaling properties, not its specific arithmetic form. A Lean proof that $A \cdot \rho^m \sigma^\ell = B$ is syntactically valid even if $\rho^m \sigma^\ell$ is the wrong quantity.

The error was \emph{semantic}, not syntactic: we proved something true, but it wasn't the thing we intended to prove.

\subsection{How It Was Caught}

\textbf{Close reading:} A human mathematician compared the Lean formulation line-by-line against Estermann's original papers and the modern expositions (Montgomery--Vaughan). The discrepancy appeared in the functional-equation structure, where the regulator appeared with explicit divisor-function semantics in the literature but only as a scaling term in the Lean formalization.

\textbf{Architectural diagnosis:} When the error was corrected in Stages 5--7, the Laurent coefficient structure changed: the exact residue at $s=1$ became $2(\gamma - \ln q)/q$ (correct), not the previous estimate (wrong). The entire downstream analysis depended on this.

\subsection{Lesson}

A formally correct proof with the wrong object is still wrong. Kernel verification certifies logical structure, not problem mapping. Fixing this required:

\begin{enumerate}[label=(\roman*)]
  \item Direct human comparison between formalization and literature
  \item Reconstruction of dependent proofs (Stages 6--8)
  \item Re-verification of the entire downstream chain
\end{enumerate}

\subsection{Structural Insight}

This error was discovered during the frontier characterization work (WP1j--WP1l), when the formalization was pushed to characterize exact cancellation mechanisms. The wrong regulator could not explain the observed cancellation behavior, forcing a re-audit of foundational machinery.

Lesson: Pushing toward real frontier work (not just replicating literature) is what forces errors into the light.

% =============================================================================
\section{Failure Mode 3: Vacuous Impossibility Claims}
\label{sec:failure3}
% =============================================================================

\subsection{The Four ``Barriers''}

The original paper claimed four ``formally proved barriers'' to RH proof strategies. Each was kernel-verified but vacuous under scrutiny.

\subsubsection{Barrier 1: Boundary-Only Interpolation is Vacuous}

\textbf{Formalized statement:}
\begin{verbatim}
theorem affineBoundaryInterpolator_half_threeHalf :
  ∀ (f g : ℝ → ℂ),
    (f (1/2) = g (1/2) ∧ f (3/2) = g (3/2))
    → (∀ s ∈ (1/2, 3/2),
       f s - g s has zero analytic content)
\end{verbatim}

\textbf{What it proves:} Any affine interpolant matching values at two boundary points contributes zero to analytic structure. True.

\textbf{Why it's vacuous:}
\begin{itemize}[leftmargin=2em]
  \item Quantifier domain: ``all affine interpolants at $\sigma = 1/2$ and $\sigma = 3/2$'' (a single construction, not a class)
  \item Claim: This ``rules out boundary-matching-only approaches''
  \item Problem: The theorem only covers affine interpolation; other boundary-matching strategies (e.g., boundary projection operators, Cauchy-integral reconstructions) are not addressed
  \item Honest restatement: ``Affine interpolation at two points carries zero analytic content'' (narrow, true, useful)
\end{itemize}

\subsubsection{Barrier 2: Singularity-Order Mismatch Cannot Be Repaired by Scalar Renormalization}

\textbf{Formalized statement:}
\begin{verbatim}
theorem no_continuous_scalar_normalization_at_zero (n : ℕ) (n_ne_0 : n ≠ 0) :
  ¬(∃ (M : ℂ → ℂ),
    ContinuousAt M 0 ∧
    M 0 ≠ 0 ∧
    ∀ s ≠ 0, s ≠ 1,
      M s * physicalNumerator n s = h15ContourAggregate n s)
\end{verbatim}

\textbf{What it proves:} No scalar $M$, continuous at $s=0$, can repair a singularity-order mismatch between two functions. True (via pole-order analysis).

\textbf{Why it's vacuous:}
\begin{itemize}[leftmargin=2em]
  \item Quantifier domain: ``all scalars $M$ continuous at $s=0$'' (restricted class)
  \item Claim: This ``rules out scalar-renormalization-only shortcuts''
  \item Problem: The theorem only covers scalars continuous at the mismatch point; singular completions (pole subtraction, residue transport, other singular-factor machinery) are explicitly excluded
  \item Honest restatement: ``A scalar continuous at the mismatch point cannot repair a singularity-order mismatch'' (true but very narrow)
\end{itemize}

\subsubsection{Barrier 3: Naive Two-Variable Bridge is Impossible}

\textbf{Formalized theorem:} \texttt{twoBoundarySlices\_force\_cornerCompatibility}

\textbf{What it claims:} A naive construction that directly matches boundary values across two dimensions is either tautological or impossible at a certified incompatible anchor.

\textbf{Why it's vacuous:}
\begin{itemize}[leftmargin=2em]
  \item The theorem is true for a specific hand-crafted two-variable object
  \item But the quantifier domain is never specified: what is ``a two-variable bridge''?
  \item The correct statement is: ``We tried this specific construction and proved it fails'' (honest)
  \item Misleading statement: ``All two-variable bridge strategies fail'' (overstated, unsupported)
\end{itemize}

\subsubsection{Barrier 4: Independent Edge Bounds Destroy Cancellation}

\textbf{Formalized theorem:} \texttt{h15BoundaryExtraction\_without\_firstOrder\_error}

\textbf{What it proves:} If you try to bound the edges of an H15 contour integral independently (without coupling), you cannot recover the required decay rate.

\textbf{Why it's vacuous:}
\begin{itemize}[leftmargin=2em]
  \item This proves that one specific strategy (independent edge bounds) fails
  \item But does NOT rule out coupled bounds, signed decompositions, or other structures
  \item The honest claim: ``Independent edge bounds are insufficient; coupled analysis is required'' (architectural discovery, not a barrier)
  \item Overstated claim: ``Edge-based approaches are impossible'' (false; coupled edge analysis works)
\end{itemize}

\subsection{The Vacuity Checklist}

To audit any kernel-verified impossibility claim, apply this checklist:

\begin{checklist}[Vacuity Audit]
\label{checklist:vacuity}
\mbox{}
\begin{enumerate}[label=\arabic*)]
  \item \textbf{State the quantified formula explicitly:}
    $$\forall x \in \text{CLASS}, \quad P(x) \to Q(x)$$
    What is the quantifier domain (CLASS)?

  \item \textbf{Define CLASS formally (not prose).}
    Replace ``all boundary-matching strategies'' with a precise set definition.
    Examples: (i) all affine functions $a + bs$ with $a, b \in \mathbb{C}$;
    (ii) all continuous functions $f: \mathbb{C} \to \mathbb{C}$ satisfying condition $X$;
    (iii) all elements of a specific Lean type or subtype.

  \item \textbf{Provide an inhabitant.}
    Exhibit an actual $x_0 \in \text{CLASS}$.
    If you cannot construct an example, the class is either empty or you don't understand its definition.

  \item \textbf{Verify non-vacuity: Does $\exists x \in \text{CLASS}$?}
    Provide a concrete instance. If no instance exists, the universal quantification is vacuous.

  \item \textbf{Kernel-verify the theorem} (standard).

  \item \textbf{Cross-check against literature.}
    Does a peer-reviewed paper make a similar claim? If not, your claim may be novel (good or bad).
    If yes, cite it precisely.

  \item \textbf{Honest restatement.}
    Rewrite the theorem claim using only the explicit quantifier domain you defined in step 2.
    Example: Instead of ``This rules out interpolation,'' say ``We proved that affine interpolation at two points carries zero analytic content.''
\end{enumerate}
\end{checklist}

\subsection{Application: Reclassifying the Four Barriers}

\begin{table}[htbp]
\centering
\small
\begin{tabular}{lllll}
\toprule
\textbf{Barrier} & \textbf{Original Claim} & \textbf{Vacuous?} & \textbf{Honest Restatement} \\
\midrule
1 & Affine interpolation impossible & YES & Affine @ 2 points carries 0 content \\
2 & Scalar fix impossible & YES & Scalar continuous @ s=0 won't work \\
3 & Two-variable bridge impossible & YES & This specific construction fails \\
4 & Independent edges insufficient & NO & Coupled analysis required (architectural) \\
\bottomrule
\end{tabular}
\caption{Vacuity audit results for the four claimed barriers.}
\label{tab:vacuity}
\end{table}

Three of four barriers were overstated. One (Barrier 4) is a genuine architectural discovery.

% =============================================================================
\section{One Real Diagnostic: The $q^2$ Majorant Loss}
\label{sec:diagnostic}
% =============================================================================

Amid the three failure modes, one substantive finding emerged: the $q^2$ majorant loss in functional-equation machinery.

\subsection{The Discovery}

In Stage 9b (vertical-strip bounds), the formalization proved:

\begin{proposition}[CutoffBudget Diagnostic]
When applying the Estermann functional equation shift (relating $D(s, a/q)$ to $D(1-s, \bar{a}/q)$), the cost is a factor of $q^2$ multiplicity in the majorant.

Consequence: Rowwise absolute summation of H15 aggregates grows unboundedly.
Numerical example:
\begin{itemize}[leftmargin=2em]
  \item At $N = 8$: normalized absolute mass = 0.58
  \item At $N = 350$: normalized absolute mass = 4.23
\end{itemize}

Therefore, absolute-majorant strategies \textbf{cannot} produce an integrable estimate as $N \to \infty$.
\end{proposition}

\subsection{Why This Is Real (Not Vacuous)}

\begin{enumerate}[label=(\roman*)]
  \item \textbf{Quantifier domain:} Specified explicitly. ``Rowwise absolute summation'' is precisely defined.

  \item \textbf{Inhabitant:} The Estermann functional-equation shift is a concrete object, not a class.

  \item \textbf{Consequence is testable:} We computed the normalized absolute mass at specific $N$ values and confirmed the growth. This is verifiable numerically.

  \item \textbf{Forces architectural change:} The $q^2$ loss is so severe that it forces a fundamental shift in proof strategy: from absolute bounds to signed cancellation machinery (Bettin--Chandee trilinear forms).
\end{enumerate}

\subsection{This Is What Sound Negative Results Look Like}

A \textbf{real barrier} (or diagnostic) in formal mathematics:
\begin{itemize}[leftmargin=2em]
  \item Names the specific obstruction (``$q^2$ multiplicity'')
  \item Quantifies its severity (growth from 0.58 to 4.23)
  \item Explains the consequence (forces strategy switch)
  \item Points to what works instead (signed cancellation)
\end{itemize}

Not a \textbf{vacuous claim}:
\begin{itemize}[leftmargin=2em]
  \item ``This class of approaches is impossible'' (undefined class)
  \item ``All boundary-matching strategies fail'' (only one tested)
  \item ``Scalar fixes are ruled out'' (covers only continuous scalars)
\end{itemize}

% =============================================================================
\section{Knowledge Graph: Methods and Deferred Evaluation}
\label{sec:kg}
% =============================================================================

During the formalization, we curated a knowledge graph of 21,942 nodes (persons, papers, theorems, proofs, lineages) and 12,859 relations.

\subsection{Current State}

The graph exists and is partially exportable (JSON, RDF). However, it lacks:

\begin{itemize}[leftmargin=2em]
  \item \textbf{Formal ontology:} No declared schema (CIDOC-CRM? OWL? Custom?)
  \item \textbf{Extraction protocol:} Unclear whether relations are automatic, manual, or hybrid
  \item \textbf{Evaluation metrics:} No precision/recall reported on relation extraction
  \item \textbf{Versioning:} No release tags or DOI assignment
  \item \textbf{License:} No explicit open-science license
\end{itemize}

\subsection{Next Steps}

The knowledge graph is a valuable artifact and warrants separate publication as a **data paper** in venues like *Journal of Open Humanities Data* or *Data Intelligence*. That paper will include:

\begin{enumerate}[label=(\roman*)]
  \item Formal schema definition (CIDOC-CRM ontology mapping)
  \item Extraction method (NLP pipeline + validation protocol)
  \item Evaluation on held-out test set (500 relations, manually validated)
  \item Zenodo DOI and CC-BY-4.0 license
  \item Interactive web explorer and bulk downloads
\end{enumerate}

Such a paper is beyond the scope of this methods paper and will be submitted separately.

% =============================================================================
\section{Lessons for LLM-Assisted Mathematical Research}
\label{sec:lessons}
% =============================================================================

\subsection{Lesson 1: Kernel Verification is Syntactic, Not Semantic}

The Lean kernel is a type-checker and proof validator. It ensures:
\begin{itemize}[leftmargin=2em]
  \item No syntax errors
  \item No logical contradictions
  \item No inference-rule violations
  \item Clear axiom tracking
\end{itemize}

It does \textbf{not} ensure:
\begin{itemize}[leftmargin=2em]
  \item The hypothesis corresponds to reality (false axioms)
  \item The formalization corresponds to the intended problem (machinery conflation)
  \item The impossibility claim covers a meaningful class (vacuity)
\end{itemize}

\subsection{Lesson 2: Vacuity Cannot Be Automated}

Detecting whether a quantifier domain is empty, whether an example actually falls in the claimed class, or whether a class is honest---these require human judgment.

Tools can help:
\begin{itemize}[leftmargin=2em]
  \item Lint checks for undefined symbols in type signatures
  \item Automated instantiation checks (does Lean's \texttt{exact} tactic find an inhabitant?)
  \item Consistency checks against literature (automated cross-reference)
\end{itemize}

But the core judgment (``Is this quantifier domain honest?'') cannot be delegated to automation.

\subsection{Lesson 3: Multi-Model Auditing Catches Single-Model Errors}

The three false axioms were discovered by Gemini 3.1 Pro's broad 1M-token context scan. Claude had written the code earlier but had not flagged the discrepancies. Consensus across six LLM families (Claude, GPT, Gemini, Mistral, DeepSeek, Kimi) confirmed each error.

Lesson: Adversarial auditing by independent models is worth the cost.

\subsection{Lesson 4: Close Reading is Non-Negotiable}

Every critical finding in this paper came from human mathematicians directly comparing:
\begin{itemize}[leftmargin=2em]
  \item Lean formalization
  \item Original published papers
  \item Modern expositions
  \item Frontier work (pushing the formalization to new discoveries)
\end{itemize}

No automated process replaced this comparison.

\subsection{Lesson 5: Pushing Toward Real Frontier Work Exposes Errors}

The regulator conflation (Failure Mode 2) was only discovered when the frontier characterization (WP1j--WP1l) required exact cancellation analysis. With the wrong regulator, the observed cancellation behavior couldn't be explained, forcing a re-audit.

Lesson: Don't just replicate literature; push toward discovery.

% =============================================================================
\section{Recommendations for Transparent LLM-Assisted Mathematics}
\label{sec:recommendations}
% =============================================================================

\subsection{For Authors}

\begin{enumerate}[label=\textbf{\arabic*.}]
  \item \textbf{Explicit axiom audit:} List all axioms (foundational + custom + hypotheses) with justification for each. Pin model versions, dates, and API IDs.

  \item \textbf{Literature cross-check:} Before publication, have a human domain expert verify every formalized theorem against its sources. Disagreement triggers investigation, not override.

  \item \textbf{Vacuity checklist (§5):} Apply this checklist to any impossibility or barrier claim. Publish the honest restatement, not the overstated version.

  \item \textbf{Close reading as mandatory stage:} Build time and resources for human mathematicians to read formalization line-by-line against literature. This is not optional.

  \item \textbf{Adversarial multi-model auditing:} When discovering errors, have multiple independent LLM families audit independently, then compare. Consensus confirms the finding.

  \item \textbf{Frontier work as error-detection:} Formalize beyond replication; push toward discovery. Errors will surface under pressure.

  \item \textbf{Version control and DOI:} Every published formalization artifact (code, graph, audit logs) gets a Zenodo DOI. Releases are tagged in git.
\end{enumerate}

\subsection{For Venues and Peer Review}

\begin{enumerate}[label=\textbf{\arabic*.}]
  \item \textbf{Require explicit axiom audit:} Do not accept a kernel-verified formalization paper without an axiom inventory and literature justification for each custom axiom.

  \item \textbf{Reject overstated barrier claims:} If a paper claims ``we proved this class of approaches is impossible,'' demand: (i) formal definition of the class, (ii) inhabitant examples, (iii) vacuity audit.

  \item \textbf{Demand vacuity certification for impossibility results:} Every ``barrier'' or ``impossibility'' theorem must include a non-vacuity proof (inhabitant + positive witness).

  \item \textbf{Cross-check against literature:} Have reviewers verify that formalized statements match their sources. Disagreement is grounds for revision or rejection.

  \item \textbf{Model versioning requirements:} Accept only papers that pin model versions (date, API ID, hash). ``Tested on Claude Sonnet'' is not reproducible; ``Claude Sonnet 5, 2026-08-04, hash=xyz'' is.
\end{enumerate}

\subsection{For the LLM/Formal Mathematics Community}

\begin{enumerate}[label=\textbf{\arabic*.}]
  \item \textbf{Develop vacuity-detection tooling:} Build Lean tactics and SMT-based checks to automatically flag:
    \begin{itemize}[leftmargin=2em]
      \item Quantified variables that appear in the conclusion but not in the hypothesis
      \item Type definitions with no declared constructors (empty types)
      \item Tactic scripts that never use a hypothesis
    \end{itemize}
    These are not definitive but can flag candidates for manual review.

  \item \textbf{Standardize axiom-tracking:} Build into Lean/Isabelle/Coq a mandatory summary output (like \texttt{\#print axioms}) that lists: (i) foundational axioms, (ii) custom axioms, (iii) explicit hypotheses in implication chains, (iv) data dependencies. This should be human-readable and appear in every build artifact.

  \item \textbf{Publish the catalog of LLM failure modes:} This paper is one case study. We need a living catalog of known failure modes, documented with examples and remediation.
\end{enumerate}

% =============================================================================
\section{Conclusion}
\label{sec:conclusion}
% =============================================================================

Machine certification of formal proofs ensures syntactic correctness but not semantic soundness. Vacuity---proving something true of an empty set, or confusing the object under study---is the dominant failure mode in LLM-assisted formalization, outweighing syntax errors.

Three failure modes appeared in our Riemann Hypothesis formalization:
\begin{enumerate}[label=(\roman*)]
  \item Legacy false axioms (caught by multi-model forensic auditing + literature review)
  \item Regulator machinery conflation (caught by close reading + frontier work pressure)
  \item Vacuous impossibility claims (caught by explicit scope auditing and the checklist in §5)
\end{enumerate}

One real architectural diagnostic emerged: the $q^2$ majorant loss forces a switch from absolute bounds to signed cancellation. This is what sound negative results look like.

For future LLM-assisted mathematical research, we recommend:
\begin{itemize}[leftmargin=2em]
  \item \checkmark{} Kernel verification is necessary, but mandatory close reading is the baseline
  \item \checkmark{} Vacuity detection requires human interpretation and cannot be fully automated
  \item \checkmark{} Multi-model auditing and adversarial testing are cost-effective
  \item \checkmark{} Pushing toward frontier work (not just replication) exposes errors
  \item \checkmark{} Explicit, honest quantifier domains prevent overstated claims
\end{itemize}

This methodology is applicable beyond the Riemann Hypothesis to any LLM-formalized mathematics.

% =============================================================================
\section*{Acknowledgments}
% =============================================================================

The author warmly thanks G\'erard Biau for proposing the original idea of exploring the Riemann Hypothesis formalization through large language models and autoformalization tools. The author also acknowledges the Sorbonne Center for Artificial Intelligence (SCAI) and Sorbonne Universit\'e for research support.

This manuscript is a working draft intended for upcoming submission on HAL.

% =============================================================================
\begin{thebibliography}{99}
% =============================================================================

\bibitem{COPE2023}
Committee on Publication Ethics.
\newblock Artificial intelligence and authorship in scholarly publishing: Guidance for editors, authors, and peer reviewers.
\newblock Position statement, 2023.

\bibitem{IwaniecKowalski2004}
H.~Iwaniec and E.~Kowalski.
\newblock \emph{Analytic Number Theory}.
\newblock American Mathematical Society Colloquium Publications, Vol. 53, 2004.

\bibitem{MontgomeryVaughan2007}
H.~L.~Montgomery and R.~C.~Vaughan.
\newblock \emph{Multiplicative Number Theory I. Classical Theory}.
\newblock Cambridge Studies in Advanced Mathematics, Vol. 97, 2007.

\bibitem{MouraUllrich2021}
L.~de~Moura and S.~Ullrich.
\newblock The {L}ean 4 theorem prover and programming language.
\newblock In \emph{Automated Deduction -- CADE 28}, LNCS 12699, pages 625--635, 2021.

\bibitem{Tenenbaum2015}
G.~Tenenbaum.
\newblock \emph{Introduction to Analytic and Probabilistic Number Theory}.
\newblock Graduate Studies in Mathematics, Vol. 163, 2015.

\end{thebibliography}

\end{document}
